\chapter{Conclusão}

O presente trabalho teve como objetivo investigar o potencial dos jogos digitais multiplayer como ferramentas para o treinamento e desenvolvimento de habilidades sociais. Para isso, foram realizadas revisões da literatura e um experimento utilizando o jogo *Among’Us*, com a participação de 15 voluntários, divididos entre um grupo experimental, composto por 9 participantes, e um grupo controle, com seis participantes. A pesquisa buscou identificar possíveis mudanças nas habilidades sociais ao longo do período experimental, combinando análises estatísticas com a avaliação das percepções dos participantes.

Os resultados quantitativos, obtidos a partir do Inventário de Habilidades Sociais (IHS-Del-Prette), demonstraram variações nos cinco fatores avaliados. No grupo experimental, foram observados aumentos médios nos fatores F1 (Enfrentamento e autoafirmação com risco), F4 (Autoexposição a desconhecidos e situações novas) e F5 (Autocontrole da agressividade). Entretanto, as análises estatísticas não identificaram diferenças significativas entre os momentos pré e pós-experimento, nem entre os grupos experimental e controle, após a aplicação da correção de Holm. Dessa forma, os resultados não permitem afirmar que a participação nas sessões tenha produzido um desenvolvimento significativo das habilidades sociais.

Apesar disso, as análises descritivas identificaram resultados favoráveis ao grupo experimental, especialmente nos fatores F1 e F4, que apresentaram as maiores diferenças positivas em relação ao grupo controle. Esses fatores estão relacionados à capacidade de se posicionar em situações interpessoais e de interagir com pessoas desconhecidas ou em situações novas, aspectos presentes nas dinâmicas sociais do jogo utilizado. Contudo, as mudanças individuais foram heterogêneas, e não foi identificado um padrão consistente entre a quantidade de sessões frequentadas e a magnitude das alterações observadas.

Por outro lado, os questionários complementares revelaram uma percepção predominantemente positiva da experiência. Todos os participantes do grupo experimental relataram ter conhecido novas pessoas e percebido algum desenvolvimento em suas habilidades sociais. Além disso, a maioria declarou sentir-se mais confortável nas interações sociais, enquanto metade afirmou ter desenvolvido novas amizades. Também foi observada uma mudança na percepção sobre o potencial dos jogos multiplayer, uma vez que, ao final do experimento, todos os participantes passaram a acreditar que esses jogos podem influenciar as habilidades sociais.

A diferença entre os resultados estatísticos e as percepções relatadas constitui um dos principais achados deste trabalho. Embora os participantes tenham identificado benefícios pessoais e sociais associados à experiência, essas percepções não foram acompanhadas por evidências estatísticas suficientes para confirmar mudanças nos fatores avaliados. Essa ausência de significância estatística pode estar relacionada, principalmente, ao tamanho reduzido da amostra, que limita o poder estatístico dos testes e dificulta a identificação de possíveis efeitos da intervenção. Isso evidencia a importância de utilizar diferentes formas de avaliação, considerando tanto os indicadores quantitativos quanto as experiências subjetivas dos participantes.

Entre as principais limitações da pesquisa, destacam-se as dificuldades relacionadas ao recrutamento e à permanência dos voluntários durante o período experimental. A dificuldade em conseguir participantes resultou em uma amostra reduzida, o que limitou o poder estatístico das análises. Além disso, a dificuldade em manter os voluntários ao longo do experimento resultou em diferentes frequências de participação, variando de uma a oito sessões, o que pode ter influenciado a intensidade da exposição às interações sociais proporcionadas pelo jogo e contribuído para a heterogeneidade das mudanças individuais observadas. Outro desafio importante foi encontrar voluntários sem experiência prévia com o jogo em questão, uma vez que sete dos oito participantes do grupo experimental já conheciam o jogo. Essa familiaridade pode ter influenciado a forma como os participantes interagiram durante as sessões, dificultando a avaliação dos efeitos da experiência em jogadores iniciantes. Por fim, fatores externos ao experimento também podem ter influenciado os resultados, especialmente no grupo controle, no qual cinco dos seis participantes relataram situações externas que poderiam afetar suas interações sociais. Dessa forma, essas limitações dificultam a atribuição das mudanças observadas exclusivamente à intervenção experimental.

Para trabalhos futuros, recomenda-se a realização de experimentos com amostras maiores, períodos de acompanhamento mais longos e maior regularidade na participação. Nesse sentido, destaca-se a possibilidade de dar continuidade à pesquisa a partir das gravações de áudio das sessões realizadas, o que permitiria investigar diretamente aspectos das interações sociais, como comunicação, cooperação, resolução de conflitos e participação nas discussões. A análise desse material, respeitando as condições de consentimento e confidencialidade da pesquisa, poderia complementar os resultados obtidos pelos questionários e contribuir para compreender possíveis diferenças entre o desenvolvimento percebido e as mudanças identificadas quantitativamente. Por fim, sugere-se investigar diferentes jogos multiplayer comerciais, considerando suas características cooperativas e competitivas, a fim de compreender melhor quais elementos dessas experiências podem estar associados ao desenvolvimento das habilidades sociais uma vez que, como relatado na revisão, a maior parte da literatura atua com jogos sérios e não jogos comerciais como foi o caso.

Conclui-se, portanto, que, embora não tenha sido possível comprovar estatisticamente a eficácia do jogo \textit{Among’Us} como ferramenta de treinamento de habilidades sociais, a experiência apresentou resultados descritivos e percepções subjetivas favoráveis à sua utilização como ambiente de interação social. Assim, o trabalho contribui para a investigação do potencial dos jogos multiplayer comerciais nesse contexto, oferecendo resultados exploratórios e uma experiência metodológica que podem orientar pesquisas futuras.